% 子文件：5.新民主主义革命理论，仅input导入
\noparindent
\section{新民主主义革命理论}

\subsection{近代中国国情}
\begin{enumerate}
    \item \textbf{认清国情}：认识解决中国革命一切基本依据；
    \item \textbf{基本国情}：半殖民地半封建社会；
    \item \textbf{主要矛盾}：帝国主义与中华民族（最主要）、封建主义与人民大众；
    \item \textbf{根本任务}：推翻三座大山，扫清复兴发展障碍；
    \item \textbf{革命性质}：资产阶级民主革命，非无产阶级社会主义革命。
\end{enumerate}

\subsection{近代中国革命时代特征}
\begin{enumerate}
    \item 十月革命后，中国民主革命转为世界无产阶级革命一部分；
    \item 五四运动标志进入新民主主义革命阶段。
\end{enumerate}

\subsection{新民主主义革命总路线}
\begin{enumerate}
    \item 1939《中国革命和中国共产党》首次提出；
    \item 1948《在晋绥干部会议上的讲话》完整概括：无产阶级领导，人民大众，反帝反封建反官僚资本主义；
    \item 作用：明确革命对象、动力、领导权、性质、前途。
\end{enumerate}

\subsection{新民主主义革命的动力}
\paragraph{人民大众组成}工人阶级（领导）、农民阶级（主力军）、城市小资产阶级、民族资产阶级。

\subsection{新民主主义革命领导力量}
\begin{enumerate}
    \item 无产阶级领导权是新旧民主革命根本区分标志；
    \item 无产阶级自身特点决定其领导地位，是革命胜利根本保证。
\end{enumerate}
\vspace{0.4em}
\noindent\sloppy\colorbox{orange!10}{
    \begin{minipage}{0.92\textwidth}
        \textbf{点拨：}实现领导权五大条件：\ding{172}联合斗争取胜；\ding{173}工农联盟为关键；\ding{174}又联合又斗争；\ding{175}革命武装作支柱；\ding{176}加强党建为根本。
    \end{minipage}
}
\vspace{0.4em}

\paragraph{核心概念对照表}
\begin{center}
\begin{tabularx}{0.95\textwidth}{|>{\centering\arraybackslash}X|>{\centering\arraybackslash}X|}
\hline
\rowcolor{orange!10} 关键命题 & 核心内容 \\
\hline
革命首要问题 & 分清敌友 \\
\hline 革命中心/理论核心 & 无产阶级领导权 \\
\hline 革命基本问题 & 农民问题 \\
\hline 革命主要内容 & 没收地主土地归农民 \\
\hline 革命实质 & 中共领导下农民革命 \\
\hline
\end{tabularx}
\end{center}

\subsection{新民主主义革命性质与前途}
\begin{enumerate}
    \item 属于新式资产阶级民主革命，区别旧民主革命；
    \item 四大新特征：世界无产革命一部分、中共领导、马列指导、前途社会主义。
\end{enumerate}

\begin{enumerate}[resume]
    \item 新民主主义革命与社会主义革命的区别——二者性质不同。
    \item 新民主主义革命与社会主义革命的联系——紧密衔接，民主革命内含社会主义因素。
    \begin{itemize}
        \item 论断：民主主义革命是社会主义革命的必要准备，社会主义革命是民主主义革命的必然趋势。
        \item 中国革命必须分“两步走”：第一步民主主义革命，第二步社会主义革命。
    \end{itemize}
    \item 认清二者关系的重要性：只有分清区别、看清联系，才能正确领导中国革命。
\end{enumerate}

\vspace{0.4em}
\paragraph{新旧民主革命与社会主义革命区分表}
\begin{center}
\begin{tabularx}{0.97\textwidth}{|>{\centering\arraybackslash}X|>{\centering\arraybackslash}X|>{\centering\arraybackslash}X|>{\centering\arraybackslash}X|>{\centering\arraybackslash}X|}
\hline
\rowcolor{orange!10} 类别 & 革命性质 & 资本政策 & 国家前途 & 二者关系 \\
\hline
新民主主义革命 & 资产阶级民主革命 & 没收官僚、保护民族资本 & 各革命阶级联合专政 & \multirow{2}{=}两步走：民主革命是准备，社会主义是必然趋势 \\
\cline{1-4}
社会主义革命 & 无产阶级革命 & 消灭私有制 & 无产阶级专政 & \\
\hline
\end{tabularx}
\end{center}

\vspace{0.4em}
\noindent\sloppy\colorbox{orange!10}{
    \begin{minipage}{0.92\textwidth}
        \textbf{点拨：}
        \ding{172} “一次革命论”（左倾无间断革命论）：只看见联系、混淆阶段区别，反帝反封建同时反对民族资产阶级，否定革命阶段性；
        \ding{173} “二次革命论”（右倾错误）：只看见区别、割裂阶段联系，放弃无产阶级对民主革命的领导权。
        \par 切勿混淆左倾“无间断革命论”与马克思主义“不断革命论”。
    \end{minipage}
}
\vspace{0.4em}

\subsection{新民主主义革命的基本纲领}
新民主主义基本纲领是总路线的具体展开，为革命指明具体奋斗目标。
\begin{enumerate}
    \item \textbf{政治纲领}：推翻帝国主义、封建主义统治，建立无产阶级领导、工农联盟为基础、各革命阶级联合专政的新民主主义共和国；属于过渡性国家形态。
    \item \textbf{经济纲领}三大内容：
    \begin{itemize}
        \item 没收封建地主土地归农民所有（新民主主义革命核心内容）；
        \item 没收官僚垄断资本归新民主主义国家所有（具备双重性质）；
        \item 保护民族工商业（由落后国情与革命性质决定，独有特色）。
    \end{itemize}
    三大经济纲领是新民主主义共和国经济建设根本方针。
\end{enumerate}
\vspace{0.4em}
\noindent\sloppy\colorbox{orange!10}{
    \begin{minipage}{0.92\textwidth}
        \textbf{点拨：}反对官僚资本主义，不是因为它属于资本主义，而是官僚资本依附帝国主义、封建地主势力，是买办封建式国家垄断资本，是国民党政权经济根基，严重阻碍中国发展。
    \end{minipage}
}
\vspace{0.4em}

\begin{enumerate}[resume]
    \item \textbf{文化纲领}：无产阶级领导、人民大众、反帝反封建的文化，即民族的、科学的、大众的文化。
    \begin{itemize}
        \item 民族的文化：反帝独立，保有民族风格特色，兼容外来进步文化；
        \item 科学的文化：反对封建迷信，坚持实事求是、理论实践统一；
        \item 大众的文化：服务工农劳动群众，面向大众、提升群众，具有民主属性。
    \end{itemize}
\end{enumerate}

\subsection{新民主主义革命道路：农村包围城市、武装夺取政权}
\subsubsection{革命道路的提出过程}
\begin{enumerate}
    \item \textbf{初步形成}：1927大革命失败后，毛泽东创建井冈山根据地，撰写《中国的红色政权为什么能够存在?》《井冈山的斗争》《星星之火，可以燎原》，提出“工农武装割据”，初步构建农村包围城市理论。
    \item \textbf{正式确立}：党的六届六中全会明确：党的任务不是长期合法斗争、先占城市后取乡村，而是走农村包围城市的相反道路。
\end{enumerate}

\subsubsection{道路形成的历史必然性}
由近代中国半殖民地半封建特殊国情决定：
\begin{enumerate}
    \item 国内无民主制度、受封建压迫；外部无民族独立、受帝国主义压迫，革命唯一斗争形式为武装斗争，核心组织是人民军队。
    \item 近代中国农业大国，农民占人口绝大多数，是无产阶级可靠同盟军、革命主力军；党必须深入农村解决土地问题，依靠农民建立革命根基，最终夺取全国政权。
\end{enumerate}